% Fill in from phase1-data/DATA_CARD.md once finalized — don't duplicate
% by hand, literally pull the numbers from there so the two documents can't
% drift out of sync.

We construct our training and evaluation corpus from [SOURCES — see
\texttt{phase1-data/DATA\_CARD.md}], totaling [N] hours across [N] speakers.
We hold out a speaker-disjoint test set of [N] utterances ([N] speakers)
for all evaluation reported in \S\ref{sec:results}; no test-set speaker
appears in any training split.

\paragraph{Licensing.} [State the Phase 1 licensing decision here —
research-only vs. permissive-only vs. both-variants track. This directly
determines what can be claimed about downstream deployability in
\S\ref{sec:ethics}.]

\paragraph{Domain.} [State whether the corpus is predominantly read speech,
scripted, or spontaneous — and whether this matches or mismatches the
conversational domain the system is evaluated on in \S\ref{sec:experiments}.
A train/eval domain mismatch is a real limitation, state it here plainly if
true rather than only in \S\ref{sec:limitations}.]

\paragraph{Pipeline validation data.} Independent of the training/eval
corpus above, we additionally used \citet{shunyalabs-yoruba} (3{,}548
utterances, unstated license) exclusively for end-to-end pipeline
smoke-testing during development --- confirming that ASR, dialogue, and TTS
components function correctly together before running the licensed
evaluation reported in \S\ref{sec:results}. We do not include this data in
any reported training or evaluation numbers, given its unclear licensing
status.
